\section{Setup}
\label{sec:methods}

\paragraph{Documents.} A seeded generator writes ``annual reports'' of
fictional networks. Each one has a title, author and date. Its sections
hold paragraphs that each plant one fact (\emph{In March, the Kadunnu
observatory recorded 120 km of cable, according to a note filed by Orka
Lolen}) plus filler sentences. There is also one bullet list, one table
(unit, budget (k\$), staff, score) and one equation. Every name is a
random syllable string, so no answer comes from pretraining. About
\NQuestionsPerDoc{} questions per document cover fact values, people,
table cells, the row with the most staff, list length and list items,
which of two facts came earlier in the year, and the equation constant.

\paragraph{Formats.} The same content tree is rendered into TXT, Markdown,
JSON (indent 2), HTML5 and LaTeX by hand-written renderers, and into DOCX,
PPTX and XLSX by python-docx, python-pptx and openpyxl. PDF is the LaTeX
compiled by pdflatex. A model can receive a binary file in two ways: as
the text a standard loader extracts (pypdf, python-docx, python-pptx,
openpyxl) or as the raw XML parts of the zip package. That gives
twelve \emph{read formats}.

\paragraph{Models.} Local, open-weight instruction models on one RTX 3060
(12\,GB): Qwen2.5-1.5B-Instruct~\cite{qwen2025qwen25},
Phi-3.5-mini-instruct~\cite{abdin2024phi3}, and
Qwen3-8B~\cite{qwen2025qwen3} in NF4 4-bit~\cite{dettmers2023qlora} with
thinking disabled. Decoding is greedy.

\paragraph{Read.} We count tokens per format under five tokenizers.
\emph{Content density} is tokens(TXT)/tokens(format). Fitting
tokens(format) $= a + b\cdot$tokens(TXT) over a size sweep (2, 4 and 8
sections) splits markup into a fixed cost $a$ and a per-content cost $b$.
Loader time is the median time to turn the file into the model's string.

\paragraph{Understand.} The prompt places the document in
\texttt{<document>} tags and appends the question. All questions about one
(document, format) share a token prefix. It is prefilled once, timed with
CUDA synchronization (the \emph{read latency}), and each question is
answered from a copy of its KV cache, as prefix caching does in serving
systems~\cite{kwon2023vllm}. A numeric answer is correct if the first
number on the first answer line matches. A text answer is correct if the
normalized gold appears in a short answer, and a comparison answer must not
also name the other entity. In the \emph{budget} condition each format is
cut to the token count of its TXT version. Formats are compared with
Markdown with exact McNemar tests~\cite{mcnemar1947} paired by question,
Holm-adjusted~\cite{holm1979}.

\paragraph{Write.} The model receives the report as plain text and one of
ten instructions, which give the thirteen deliverables of
Table~\ref{tab:write}. Each reply's first fenced block
is built and validated:
\begin{itemize}
\item JSON must parse.
\item LaTeX must compile with pdflatex.
\item HTML and Markdown are turned into PDF with xhtml2pdf (via
  python-markdown for Markdown).
\item python-docx, python-pptx and openpyxl code runs in a subprocess with
  a 60\,s timeout and must save a file that opens.
\item Raw \texttt{word/document.xml} is packed into a blank .docx that
  python-docx must open.
\item CSV must have at least 2 rows of equal width.
\end{itemize}
A failed build is returned to the model with the error, up to
\MaxAttempts{} attempts. The cost of a deliverable is its \emph{time per
valid file}: all generation and build time, failures included, divided by
the number of valid files. The \textbf{format tax} is that time divided by
the same quantity for Markdown. Fidelity is the share of planted strings
that the built file, read back by the loaders, still contains.
Spreadsheet routes write only the table and are scored on its cells.
